\medskip

%\begin{minipage}{0.58\linewidth}
%Sur la figure ci-contre :
%
%\begin{itemize}[label=$\bullet~~$]
%\item ABCD est un carré de côté 9 cm; 
%\item les segments de même longueur sont codés.
%\end{itemize}

\begin{enumerate}
\item 
	\begin{enumerate}
		\item %Le polygone IJKLMNOP est-il régulier, c'est-à-dire a-t-il
%tous ses côtés de même longueur ?
Les côtés obliques du polygone (c'est un octogone) IJKLMNOP sont tous les hypoténuses de triangles rectangles isocèles dont les côtés de l'angle droit mesurent $\dfrac93 = 3$~(cm).

Par exemple dans le triangle rectangle en A, AIP, on a 

\[\text{AI}^2 + \text{AP}^2 = \text{IP}^2\:\text{ou } 3^2 + 3^2 = \text{IP}^2\]

On a donc IP$^2 = 9 + 9 = 18$ ; on en déduit que 

IP $ = \sqrt{18} = \sqrt{9 \times 2} = \sqrt{9} \times \sqrt 2 = 3\sqrt 2\ne 3$.

L'octogone a 4 côtés de longueurs 3 et 4 côtés de longueur $3\sqrt{2} \approx 4,24$. : il n'estpas régulier

\textbf{Justifier la réponse.}
		\item %Justifier que l’aire de la surface IJKLMNOP grisée sur
%la figure ci-contre est égale à 63~cm$^2$.
l'aire grisée est égale à différence entre l'aire du carré ABCD et les quatre triangles rectangles isocèles en blanc :

$\bullet~$ aire du carré ABCD : $9^2 = 81$~(cm$^2$)

$\bullet~$ aire des quatre triangles : $4 \times \dfrac{3 \times 3}{2} = 2 \times 9 = 18$~(cm$^2$).

L'aire cherchée est égale à :

\[81 - 18 = 63~(\text{cm}^2.\]

	\end{enumerate}
%\end{enumerate}
%\end{minipage}\hfill
%\begin{minipage}{0.38\linewidth}
%\psset{unit=0.9cm}
%\begin{pspicture}(0,-1)(5.8,5.8)
%%\psgrid
%\psframe(0.2,0.2)(5.6,5.6)
%\pspolygon[fillstyle=solid,fillcolor=lightgray](2,5.6)(3.8,5.6)(5.6,3.8)(5.6,2)(3.8,0.2)(2,0.2)(0.2,2)(0.2,3.8)%IJKLMNOP
%\def\dblv{\psline(-0.05,0.1)(-0.05,-0.1)\psline(0.05,0.1)(0.05,-0.1)}
%\rput(2.9,5.6){\dblv}\rput(1,5.6){\dblv}\rput(4.7,5.6){\dblv}
%\rput(2.9,0.2){\dblv}\rput(1,0.2){\dblv}\rput(4.7,0.2){\dblv}
%\def\dblh{\psline(-0.1,0.05)(0.1,0.05)\psline(-0.1,-0.05)(0.1,-0.05)}
%\rput(0.2,4.7){\dblh}\rput(0.2,2.9){\dblh}\rput(0.2,1.1){\dblh}
%\rput(5.6,4.7){\dblh}\rput(5.6,2.9){\dblh}\rput(5.6,1.1){\dblh}
%\uput[ul](0.2,5.6){A} \uput[ur](5.6,5.6){B} \uput[dr](5.6,0.2){C} \uput[dl](0.2,0.2){D}
%\uput[u](2,5.6){I} \uput[u](3.8,5.6){J} \uput[r](5.6,3.8){K} \uput[r](5.6,2){L}
%\uput[d](3.8,0.2){M} \uput[d](2,0.2){N} \uput[l](0.2,2){O} \uput[l](0.2,3.8){P}
%\end{pspicture}
%\end{minipage}

%\begin{minipage}{0.58\linewidth}
%\begin{enumerate}%[start=2]
\item %Les diagonales du carré ABCD se coupent en S.

%\smallskip

%On a tracé le cercle de centre S et de diamètre 9 cm.
	\begin{enumerate}
		\item %Déterminer l'aire du disque de centre S et de
%diamètre 9 cm.
Le disque a un rayon de $\dfrac92 = 4,5$ ; son aire est donc égale à $\pi R^2 = \pi \times 4,5^2 = 20,25\pi$.
		\item %Montrer que la différence entre l'aire du polygone
%IJKLMNOP et l'aire du disque représente moins de
%1\,\% de l’aire du disque.
Comme (calculette) $20,25\pi \approx 63,62$~(cm$^2$) c'est le disque qui a la plus grande aire.

On calcule le pourcentage de la différence entre l'aire du polygone
-IJKLMNOP et l'aire du disque par rapport à l'aire du disque en calculant:

$\dfrac{\text{aire du disque} - \text{aire du polygone}}{\text{aire du disque}}\times 100$, soit 

$\dfrac{20,25\pi - 63}{20,25\pi} \times 100 \approx 0,97 < 1$.

Ce pourcent est bien inférieur à 1.
	\end{enumerate}
\end{enumerate}
%\end{minipage}\hfill
%\begin{minipage}{0.38\linewidth}
%\psset{unit=0.9cm}
%\begin{pspicture}(0,0)(5.8,5.8)
%%\psgrid
%\psframe(0.2,0.2)(5.6,5.6)
%\pspolygon(2,5.6)(3.8,5.6)(5.6,3.8)(5.6,2)(3.8,0.2)(2,0.2)(0.2,2)(0.2,3.8)%IJKLMNOP
%\pscircle[fillstyle=solid,fillcolor=lightgray](2.9,2.9){2.7}
%\def\dblv{\psline(-0.05,0.1)(-0.05,-0.1)\psline(0.05,0.1)(0.05,-0.1)}
%\rput(2.9,5.6){\dblv}\rput(1,5.6){\dblv}\rput(4.7,5.6){\dblv}
%\rput(2.9,0.2){\dblv}\rput(1,0.2){\dblv}\rput(4.7,0.2){\dblv}
%\def\dblh{\psline(-0.1,0.05)(0.1,0.05)\psline(-0.1,-0.05)(0.1,-0.05)}
%\rput(0.2,4.7){\dblh}\rput(0.2,2.9){\dblh}\rput(0.2,1.1){\dblh}
%\rput(5.6,4.7){\dblh}\rput(5.6,2.9){\dblh}\rput(5.6,1.1){\dblh}
%\uput[ul](0.2,5.6){A} \uput[ur](5.6,5.6){B} \uput[dr](5.6,0.2){C} \uput[dl](0.2,0.2){D}
%\uput[u](2,5.6){I} \uput[u](3.8,5.6){J} \uput[r](5.6,3.8){K} \uput[r](5.6,2){L}
%\uput[d](3.8,0.2){M} \uput[d](2,0.2){N} \uput[l](0.2,2){O} \uput[l](0.2,3.8){P}\uput[l](2.9,2.9){S}
%\psline[linestyle=dashed,linewidth=1.5pt](0.2,5.6)(5.6,0.2)
%\psline[linestyle=dashed,linewidth=1.5pt](5.6,5.6)(0.2,0.2)
%\pspolygon(2,5.6)(3.8,5.6)(5.6,3.8)(5.6,2)(3.8,0.2)(2,0.2)(0.2,2)(0.2,3.8)%IJKLMNOP
%\end{pspicture}
%\end{minipage}

\emph{Remarque} : quand on ne connaissait pas la formule donnant l'aire du disque on cherchait à approcher cette aire par celle d'un polygone dont le contour voisinait le cercle.

On voit avec cet exercice que l'octogone ci-dessus donne une aire (63) du disque avec une erreur inférieure à $1\,\%$.
